%! AP Statistics — Unit 4, Lesson 4.5 — Guided Notes KEY
\documentclass[10pt]{article}
\usepackage{apstats-article}
\usepackage{apstats-key}

\begin{document}

\pageheader{Unit 4, Lesson 4.5}{Guided Notes --- KEY}
\namedateperiod

\vspace{0.12in}

\begin{objectivebox}
\small By the end of this lesson, I will be able to\dots\\[4pt]
\ans{Calculate conditional probabilities using the formula $P(A|B) = P(A \cap B)/P(B)$.}\\[1pt]
\ans{Read a two-way table to identify the joint probability and the conditioning event total.}\\[1pt]
\ans{Apply the multiplication rule $P(A \cap B) = P(A) \cdot P(B|A)$ to find a joint probability.}
\end{objectivebox}

\vspace{0.1in}

\begin{vocabbox}
\termblanklong{Conditional probability $P(A|B)$} \ans{The probability that event $A$ will occur given that event $B$ has already occurred; $P(A|B) = P(A \cap B)/P(B)$.}
\termblanklong{Conditioning event} \ans{The event in the denominator of the conditional probability formula; the ``given'' event that restricts the sample space.}
\termblanklong{Multiplication rule} \ans{$P(A \cap B) = P(A) \cdot P(B|A)$; used to find joint probability from marginal and conditional probabilities.}
\termblanklong{Marginal probability} \ans{The probability of a single event, found from a row or column total in a two-way table.}
\end{vocabbox}

\vspace{0.1in}

\begin{hookbox}
\small
\textbf{Card scenario:}
\begin{enumerate}[leftmargin=1.5em, itemsep=4pt]
  \item \ans{$P(\text{King}) = 4/52 = 1/13 \approx 0.077$}
  \item \ans{$P(\text{King}|\text{face card}) = 4/12 = 1/3 \approx 0.333$}
  \item \ans{12 face cards; 4 of them are Kings.}
\end{enumerate}
\ans{The probability increased because knowing the card is a face card restricts the sample space from 52 cards to only the 12 face cards.}
\end{hookbox}

\vspace{0.1in}

\begin{notesbox}{Conditional Probability (VAR-4.D.1)}
\small
\textbf{Definition:} The \textbf{conditional probability} $P(A|B)$ is the probability
that event $A$ will occur \ans{given that} event $B$ \ans{has occurred}.

\vspace{8pt}
\textbf{Formula:}
\[
  P(A|B) = \frac{P(A \cap B)}{P(B)}, \quad \text{provided } P(B) \ans{> 0}.
\]

\vspace{6pt}
\textbf{Interpretation:} Knowing $B$ occurred \ans{restricts} the sample space to only
outcomes in $B$.

\vspace{10pt}
\textbf{From a two-way table:}
\begin{enumerate}[leftmargin=1.5em, itemsep=4pt]
  \item Find $P(A \cap B)$ = \ans{joint cell count} $\div$ grand total
  \item Find $P(B)$ = \ans{conditioning row or column total} $\div$ grand total
  \item Compute $P(A|B)$ = \ans{joint cell count} $\div$ \ans{conditioning total} (the grand total cancels!)
\end{enumerate}

\vspace{8pt}
\textbf{Short-cut for tables:} $P(A|B) = \dfrac{\text{count in both }A\text{ and }B}{\ans{\text{count in }B}}$
\end{notesbox}

\vspace{0.1in}

\begin{notesbox}{Worked Example --- Two-Way Table}
\small
\vspace{4pt}
\renewcommand{\arraystretch}{1.4}
\begin{tabularx}{\linewidth}{@{} l X X r @{}}
 & \textbf{Trains Daily} & \textbf{Trains Weekly} & \textbf{Total} \\
\hline
\textbf{Runs}     & 24 & 16 & 40 \\
\textbf{Swims}    & 18 & 22 & 40 \\
\textbf{Cycles}   & 20 & 20 & 40 \\
\hline
\textbf{Total}    & 62 & 58 & 120 \\
\end{tabularx}

\vspace{8pt}
\textbf{Find $P(\text{Runs} | \text{Trains Daily})$:}
\begin{itemize}[itemsep=3pt]
  \item $P(\text{Runs} \cap \text{Trains Daily}) = $ \ans{$24/120$}
  \item $P(\text{Trains Daily}) = $ \ans{$62/120$}
  \item $P(\text{Runs} | \text{Trains Daily}) = \dfrac{\ans{24/120}}{\ans{62/120}} = \dfrac{24}{62} \approx \ans{0.387}$
\end{itemize}

\vspace{6pt}
\textbf{Interpret in context:} Among athletes who train daily,
\ans{38.7}\% run.
\end{notesbox}

\vspace{0.1in}

\begin{notesbox}{Multiplication Rule (VAR-4.D.2)}
\small
\textbf{Rearranging the formula} gives the \textbf{multiplication rule}:
\[
  P(A \cap B) = P(A) \cdot P(B|A) = P(B) \cdot P(A|B)
\]

\vspace{6pt}
\textbf{Use this when:} you know $P(A)$ and $P(B|A)$ and need to find \ans{$P(A \cap B)$}.

\vspace{8pt}
\textbf{Example:}
\[
  P(\text{Runs} \cap \text{Trains Daily}) = \textcolor{keyred}{\tfrac{62}{120}} \times \textcolor{keyred}{\tfrac{24}{62}} = \textcolor{keyred}{\tfrac{24}{120}} = \textcolor{keyred}{0.2}
\]
\end{notesbox}

\vspace{0.1in}

\begin{practicebox}
\small
\vspace{6pt}
\begin{enumerate}[leftmargin=1.5em, itemsep=10pt]
  \item \ans{$18 \div 40 = 0.45$} (18 swimmers train daily; 40 swimmers total)

  \item \ans{$20 \div 58 \approx 0.345$} (20 cyclists train weekly; 58 weekly trainers total)

  \item \ans{$(58/120) \cdot (20/58) = 20/120 \approx 0.167$. Yes, matches the cell count 20.}

  \item \ans{$P(\text{Runs}|\text{Trains Daily}) = 24/62 \approx 0.387$; $P(\text{Trains Daily}|\text{Runs}) = 24/40 = 0.6$. They are not equal --- the order of conditioning matters.}
\end{enumerate}
\end{practicebox}

\end{document}
